\section{12. Gobernanza de datos}
\label{sec:gobernanza}


\begin{consigna}{red}
En esta sección se debe analizar de manera integral el marco normativo y ético asociado al uso de los datos y al desarrollo de soluciones basadas en inteligencia artificial.

El análisis se divide en dos partes:

\begin{itemize}
  \item Cumplimiento normativo.
  \item Ética en el uso de inteligencia artificial.
\end{itemize}

\subsection{Cumplimiento normativo}
Este análisis es clave para garantizar el cumplimiento normativo y evitar conflictos legales durante el desarrollo y publicación del proyecto.

En esta subsección se debe identificar si los datos utilizados en el proyecto están sujetos a normativas de protección de datos y privacidad, y en qué condiciones pueden emplearse.

Aspectos a considerar:

\begin{itemize}
  \item Determinar si los datos están regulados por normativas como el GDPR, la Ley 25.326 de Protección de Datos Personales (Argentina), la HIPAA u otras, según la jurisdicción o la temática del proyecto.
  \item Aclarar si las fuentes de datos son propias, de acceso público o licenciadas. En este último caso, explicitar la licencia correspondiente y sus condiciones de uso. Realizar este mismo análisis para las herramientas y, en caso de corresponder, para los modelos preentrenados a emplear.
  \item Analizar la viabilidad legal del proyecto, considerando los mecanismos necesarios para garantizar el cumplimiento normativo y la trazabilidad de los datos.
\end{itemize}

\subsection{Ética en el uso de inteligencia artificial}

En esta subsección se debe reflexionar sobre los aspectos éticos vinculados al uso de datos y algoritmos de inteligencia artificial. Se espera una evaluación de la integridad del sistema y su impacto social.

Aspectos a considerar:

\begin{itemize}
	\item Identificar posibles sesgos en los datos o modelos (ideológicos, culturales, de género, geográficos, etc.) y analizar sus consecuencias (por ejemplo: discriminación, manipulación, pérdida de privacidad o desinformación).
	\item Analizar los posibles riesgos en términos de impacto social del uso indebido de la solución a desarrollar.
	\item En caso de corresponder, proponer medidas de mitigación y mecanismos de confianza (auditorías, documentación, trazabilidad, revisión humana, etc.).
	\item Finalmente, elaborar una reflexión general sobre la ética del proyecto, considerando tanto la equidad y transparencia del sistema como su impacto potencial en la sociedad.

\end{itemize}
\end{consigna}
